% =============================================================================
\chapter{Robustness}
\label{ch:robustness}
% =============================================================================

% TODO: Structure
%
% \section{Alternative Volatility Windows}
% - 63-day and 180-day post-filing windows
% - Check whether divergence effect is short-lived or persistent
%
% \section{Stock Returns as Alternative Dependent Variable}
% - Re-estimate all three models with post-filing log returns as DV
% - Weaker signal expected (low cross-sectional R²); included for completeness
%
% \section{Alternative Divergence Constructions}
% - Swap ROA for operating\_roa or ocf\_to\_assets in ΔROA
% - Use binary DivDummy instead of continuous divergence
% - Residual-based divergence (regress ΔSent on ΔROA, take residual)
%
% \section{Alternative Textual Measures}
% - Different dictionaries (Henry 2008 vs.\ Loughran--McDonald)
% - Tf-idf weighting instead of raw word counts
%
% \section{Subsample Analysis}
% - Pre/post-2020 (COVID, supply-chain disruptions, inflation)
% - High-tech vs.\ traditional industries
%
% \section{Sensitivity to Sample Construction}
% - Excluding firms with fewer than 8 years of data
% - Effect of fiscal year-end alignment
% - Winsorisation at 2.5\%/97.5\%

% ─────────────────────────────────────────────────────────────────────────────
\section{Alternative post-filing windows}
\label{sec:robustness_horizons}
% ─────────────────────────────────────────────────────────────────────────────

The main specifications measure post-filing volatility over a
365-day window. One year matches the annual frequency of the panel,
but it has two drawbacks. First, the window is long enough that the
next 10-K is filed before it ends, so the second half of the window
already contains information from the next disclosure. Second, prior
work on text in 10-Ks consistently finds that any predictive power
shows up in the first weeks or months after the filing and fades
afterwards \parencite{loughran2011}. A one-year window
mixes the short signal with many months of unrelated noise.

To check whether the main results depend on this choice, the
divergence model (equation~\ref{eq:divergence_model}) is re-estimated
with the dependent variable replaced by volatility measured over
30, 90, and 180~days after the filing. The lagged dependent variable
is the same firm's volatility computed over the previous filing
year's window of the same length. All other controls, fixed effects,
and the firm-clustered standard errors are kept the same.
Table~\ref{tab:horizons} shows all four horizons side by side; the
365-day column reproduces the divergence model from
Chapter~\ref{ch:results}.

\begin{table}[htbp]
\centering
\caption{Post-Filing Volatility at Alternative Horizons}
\label{tab:horizons}
\begin{tabular}{lcccc}
\toprule
 & (1) 30d & (2) 90d & (3) 180d & (4) 365d \\
\midrule
Lagged Volatility & $+0.0888$*** & $+0.1921$*** & $+0.3150$*** & $+0.5361$*** \\
 & $(+3.97)$ & $(+8.20)$ & $(+14.66)$ & $(+22.69)$ \\[4pt]
$\Delta$Sentiment & $+14.0930$*** & $+11.3194$*** & $+8.6803$** & $+6.5744$*** \\
 & $(+2.94)$ & $(+2.89)$ & $(+2.57)$ & $(+2.66)$ \\[4pt]
$\Delta$Risk & $+0.1119$ & $-0.5747$ & $-0.6601$ & $-0.0707$ \\
 & $(+0.06)$ & $(-0.45)$ & $(-0.66)$ & $(-0.11)$ \\[4pt]
TextSim & $+0.0501$ & $-0.0045$ & $+0.0462$ & $+0.0620$* \\
 & $(+0.68)$ & $(-0.07)$ & $(+0.99)$ & $(+1.83)$ \\[4pt]
Divergence & $-14.6169$*** & $-11.0778$*** & $-7.9999$** & $-4.8501$* \\
 & $(-3.03)$ & $(-2.79)$ & $(-2.31)$ & $(-1.86)$ \\[4pt]
\midrule
Industry FE & Yes & Yes & Yes & Yes \\
Year FE & Yes & Yes & Yes & Yes \\
Observations & 5,252 & 5,252 & 5,252 & 5,473 \\
Firms & 510 & 510 & 510 & 531 \\
Adj.\ $R^2$ & 0.447 & 0.500 & 0.535 & 0.645 \\
\bottomrule
\end{tabular}
\begin{tablenotes}
\item \textit{Note:} Each column re-estimates the divergence specification (equation~\ref{eq:divergence_model}) with the dependent variable replaced by the annualised standard deviation of daily log returns over a shorter post-filing window. The lagged dependent variable is correspondingly the same firm's volatility measured over the prior filing year's window of the same length. All specifications include two-digit SIC industry and fiscal-year fixed effects; $t$-statistics in parentheses use firm-clustered standard errors. ***, **, and * denote significance at the 1\%, 5\%, and 10\% levels, respectively.
\end{tablenotes}
\end{table}

The pattern is clear and goes in one direction. The divergence
coefficient is largest at 30~days ($\delta = -14.62$, $t = -3.03$,
$p = 0.002$), shrinks to $-11.08$ at 90~days, falls further to
$-8.00$ at 180~days, and reaches $-4.85$ at one year, where it
becomes only marginally significant. The size of the effect drops
roughly by a factor of three from the shortest to the longest window,
and the $t$-statistic moves in the same direction. The same
monotonic decay is visible for $\Delta$Sentiment, while
$\Delta$Risk and TextSim remain small at every horizon.

Two conclusions follow. First, the headline divergence result is not
an artefact of the one-year window: the effect is in fact stronger
at shorter horizons and is highly significant at every sub-annual
window. The marginal $p$-value in the annual specification reported
in Chapter~\ref{ch:results} reflects dilution by months of unrelated
post-filing noise, not the absence of a signal. H2 is therefore
supported once the dependent variable is matched to the horizon over
which the textual signal actually operates. Second, the decay
matches the timing patterns documented for textual measures in
10-Ks by \textcite{loughran2011}, which adds external validity to
the divergence construction used here.
